\scanpage{87}
% loc 1–2
(l'invariant \uline{anharmonique} de $k[u]$ et $k[v]$ dans alg. d'Azumaya)
% loc 3
\[ \deleted{\unsure{\text{de}}}\quad \rho=f(b,c,b',c';t)\deleted{\unsure{a}}=\frac{P(b,c,b',c',t)}{Q(\text{---})} \]
% loc 4–7
telle que la condition néc. et suffisante pour que \deleted{$M$}$(b,c,b',c',t)$ \deleted{soit alg. d'Azumaya}\inserted{se réalise loc. fpqc par alg. d'Azumaya}, \deleted{c'est que} (en supposant déjà $\delta\delta'$ inversible), c'est que $\rho\inserted{\text{partout}}\ne2$,
% loc 8
\[ \text{i.e.}\qquad P-2Q\inserted{\text{partout}}\ne2 \]
% loc 9–18
En fait, quand ils sont réalisés dans alg. d'Azumaya, \deleted{ils définissent} alors $A_1,A_2$ définissent \deleted{deux div} chacun un diviseur de degré 2\deleted{, et} sur \deleted{\unsure{l'id}} le schéma projectif associé à $A$, et si leurs supports sont disjoints, il y a un inv. anharmonique \deleted{``\unsure{défini|fois}''}$\lambda$ (qui n'est défini que localement \inserted{sur $\operatorname{Spec}k$} et mod involution $\lambda\mapsto\lambda^{-1}$) et l'invariant (défini sur $k$)
% loc 19
\[ \bigl(\qquad \rho=\lambda+\deleted{\frac{1}{\lambda}}\lambda^{-1}\qquad \text{d'où}\quad \lambda^2-\rho\lambda+1=0 \]
% loc 20–21
(et comme $\lambda$ partout $\ne0,1$, on doit avoir $\rho$ partout $\ne2\ \ldots$
% archive sheet number: 86
